% --- Asegúrate de tener esto en tu preámbulo ---


% --- Definir colores (opcional, pero mejora la legibilidad) ---
\definecolor{codegreen}{rgb}{0,0.6,0}
\definecolor{codegray}{rgb}{0.5,0.5,0.5}
\definecolor{codepurple}{rgb}{0.58,0,0.82}
\definecolor{backcolour}{rgb}{0.98,0.98,0.98} % Un fondo ligeramente gris

% --- Configuración global para lstlisting (puedes ajustarla) ---
\lstdefinestyle{mystyle}{
    backgroundcolor=\color{backcolour},
    commentstyle=\color{codegreen},
    keywordstyle=\color{magenta},
    numberstyle=\tiny\color{codegray},
    stringstyle=\color{codepurple},
    basicstyle=\ttfamily\footnotesize, % Fuente más pequeña
    breakatwhitespace=false,
    breaklines=true, % Permitir saltos de línea automáticos
    captionpos=b,
    keepspaces=true,
    numbers=left,
    numbersep=5pt,
    showspaces=false,
    showstringspaces=false,
    showtabs=false,
    tabsize=2,
    frame=single, % Mantener el marco
    framerule=0.4pt,
    rulecolor=\color{black!30}, % Color del marco más suave
    xleftmargin=15pt, % Margen izquierdo
    framexleftmargin=15pt
}
\lstset{style=mystyle} % Aplicar el estilo globalmente

% --- Inicio de la Sección ---

\section{Agregar una nueva metaheurística}
\label{sec:agregar_mh}

El diseño modular del solver facilita la incorporación de nuevas metaheurísticas bio-inspiradas o de otro tipo. Para añadir un nuevo algoritmo, sigue estos pasos:

\subsection{Paso 1: Implementación de la lógica del algoritmo}

\begin{enumerate}
    \item \textbf{Crear el archivo Python:} \\
    Navega hasta el directorio \texttt{src/metaheuristics/Codes/}. Crea un nuevo archivo Python para tu metaheurística (ej. \texttt{NAE.py}).

    \item \textbf{Definir la función de iteración:} \\
    Dentro de tu nuevo archivo, define la función principal que ejecutará una iteración del algoritmo (ej. \texttt{iterarNAE(...)}).

    \item \textbf{Establecer la firma estándar de la función:} \\
    Es \textbf{crucial} que tu función utilice los nombres de parámetros estandarizados que espera la función \texttt{iterate\_population}. La firma general podría ser:

    % --- Bloque de código Python mejor formateado ---
    \begin{lstlisting}[language=Python, caption={Ejemplo firma estándar (NAE) Mejorada}, label={lst:firma_nae_mejorada}]
import numpy as np
# otras importaciones...

def iterarNAE(maxIter, iter, dim, population, fitness, best, fo=None, lb=None, ub=None, lb0=None, ub0=None, vel=None, pBest=None, objective_type='MIN')

    """
    Realiza una iteracion del Nuevo Algoritmo de Ejemplo (NAE).

    Args:
        maxIter (int): Maximo de iteraciones.
        iter (int): Iteracion actual.
        dim (int): Dimension del problema.
        population (np.ndarray): Poblacion actual.
        fitness (np.ndarray): Fitness de la poblacion actual.
        best (np.ndarray): Mejor solucion global hasta ahora.
        fo (callable, optional): Funcion objetivo. Defaults to None.
        lb0 (float, optional): Limite inferior escalar. Defaults to None.
        ub0 (float, optional): Limite inferior escalar. Defaults to None.
        ...

    Puede tener mas argumentos segun la metaheuristica a implementar.

    Returns:
        np.ndarray or tuple: Nueva poblacion, o tupla (pob, vel/mejoras).
    """
    N = population.shape[0]
    new_population = population.copy()

    # -----------------------------------
    #    LOGICA DE TU METAHEURISTICA NAE
    # -----------------------------------

    # Ejemplo placeholder:
    # for i in range(N):
    #     r = np.random.rand(dim)
    #     new_population[i] = population[i] + r * (best - population[i])

    # --- Valor de Retorno ---
    return new_population
    # O: return new_population, new_vel
    # O: return new_population, posibles_mejoras
    \end{lstlisting}
    
    % --- Fin bloque de código ---

    \textbf{Importante:}
    \begin{itemize}
        \item No todos los parámetros son necesarios para cada algoritmo.
        \item Usa los nombres estándar exactos (\texttt{iter}, \texttt{best}, \texttt{fo}, etc.).
        \item Devuelve la nueva población (y opcionalmente velocidad o mejoras).
        \item Para SCP/USCP, la función opera en continuo; la binarización/reparación es posterior.
    \end{itemize}
\end{enumerate}

\subsection{Paso 2: Registro en el sistema}

Registra la nueva función:

\begin{enumerate}
    \item \textbf{Editar \texttt{src/metaheuristics/imports.py}:}
    \item \textbf{Añadir importación:}
        \begin{lstlisting}[language=Python]
from .Codes.NAE import iterarNAE # Ajusta NAE
        \end{lstlisting}
    \item \textbf{Añadir a \texttt{metaheuristics}:}
        \begin{lstlisting}[language=Python]
metaheuristics = {
    # ...
    "NAE": iterarNAE, # Nueva entrada
    # ...
}
        \end{lstlisting}
    \item \textbf{Añadir a \texttt{MH\_ARG\_MAP}:}
        \begin{lstlisting}[language=Python]
MH_ARG_MAP = {
    # ...
    # Lista los strings de params que tu func necesita:
    'NAE': ('maxIter', 'iter', 'dim', 'population', 'best'),
    # ...
}
        \end{lstlisting}
        ¡Asegúrate de que los nombres en la tupla coincidan \textbf{exactamente} con los parámetros que tu función \texttt{iterarNAE} usa!
\end{enumerate}

\subsection{Paso 3: Habilitar en experimentos}

Para usar la nueva MH:

\begin{enumerate}
    \item \textbf{Editar \texttt{src/util/json/experiments\_config.json}:}
    \item \textbf{Añadir a \texttt{"mhs"}:}
        \begin{lstlisting}
  // ...
  "mhs": ["SBOA", "PSO", "NAE"], // <--- Agrega aqui
  // ...
}
        \end{lstlisting}
\end{enumerate}

\subsection{Paso 4: Consideraciones adicionales (opcional)}

\begin{itemize}
    \item \textbf{Versiones específicas (BEN vs SCP/USCP):} Si necesitas lógica muy diferente, crea funciones separadas (ej. \texttt{iterarNAE\_Ben}, \texttt{iterarNAE\_Scp}) y regístralas con claves distintas ("NAE\_BEN", "NAE\_SCP"). Tengan en cuenta que idealmente debe solo haber una lógica para todos los problemas.
    \item \textbf{Algoritmos binarios (GA):} Pueden requerir manejo especial si no siguen el flujo estándar continuo + binarización.
    \item \textbf{Pruebas}: Verifica tu nueva MH ejecutándola en problemas de prueba.
\end{itemize}
Siguiendo estos pasos, podrás integrar nuevas metaheurísticas de manera organizada.

% --- Fin de la Sección ---